\documentclass[handout]{beamer}
\mode<presentation>{\usetheme{Warsaw}\usecolortheme{crane}}
\usepackage{centernot}
\usepackage{graphicx}

\usepackage[utf8]{inputenc}
\usepackage[english]{babel}
\usepackage{lmodern}
\usepackage[T1]{fontenc}
\usepackage[babel=true]{microtype}

%% preamble
\title{UNIX Weapons School---Prelude}
\date{}
\author{Nick Black for the\\
Georgia Institute of Technology
}

\begin{document}

\begin{frame}
\titlepage
\begin{center}
\includegraphics[scale=0.33]{uws.png}
\end{center}
\end{frame}

\begin{frame}{Three review questions}
\begin{itemize}
\item \textbf{Q1.} Can adding a NOP (no-operation) instruction accelerate a program?
\item \textbf{Q2.} Will a program always incur fewer cache misses on a fully-associative
	cache than a directly-mapped cache of the same size?
\item \textbf{Q3.} A machine has two symmetric unicore physical processors with HyperThreading
	(2-way SMT). Are there threads such that a single-processor
	schedule can beat using both?
\end{itemize}
\end{frame}

\begin{frame}{Review Question I}
\begin{itemize}
\item \textbf{Q1.} Can adding a NOP (no-operation) instruction accelerate a program?
\item \textbf{A1.} Yes. An example might be when it causes a branch target to 
 be aligned.
\item \textbf{$NB1_1$.} This is generally the compiler's job.
\end{itemize}
\end{frame}

\begin{frame}{Review Question II}
\begin{itemize}
\item \textbf{Q2.} Will a program always incur fewer cache misses on a fully-associative
	cache than a directly-mapped cache of the same size?
\item \textbf{A2.} No. A fully-associative cache will yield fewer hits when
  data that \textit{is not} reused evicts data which \textit{is} reused, and
  would otherwise have gone unevicted.
\item \textbf{$NB2_1$.} This implies that a cache can have a
  negative number of conflict misses!
\end{itemize}
\end{frame}

\begin{frame}{Review Question III}
\begin{itemize}
\item \textbf{Q3.} A machine has two symmetric unicore physical processors with HyperThreading
	(2-way SMT). Are there threads such that a single-processor
	schedule can beat using both?
\item \textbf{A3.} Yes. The two logical cores share all cache levels. If e.g.
	the MESI protocol is in use, and writes to shared data are frequent,
	and the dataset fits within one processor's cache, the reduction in
	cycle capacity might be more than offset by fewer RFOs. 
\end{itemize}
\end{frame}


\end{document}
